\chapter{Puissances et racines carrées}
\label{exercices:04}
\section*{Exercice 4.1 -- Les exposants ont un sens}
Calculer exactement $2^{-3}$ et $(3^2)^3/3^4$. Écrire en notation scientifique $(8\times10^7)/(2\times10^{-3})$.

\section*{Exercice 4.2 -- Racines et géométrie}
Un carré a une aire de $98\unite{cm^2}$. Donner son côté exact, le simplifier si possible, puis l'arrondir au millimètre. Calculer aussi $\sqrt{(-8)^2}$.

\section*{Exercice 4.3 -- Doublements}
Dans un modèle, $150$ cellules doublent toutes les $30$ minutes. Combien en prévoit-on après $4$ heures ? Pourquoi ne faut-il pas utiliser l'exposant $4$ ?

\section*{Exercice 4.4 -- Réfuter puis expliquer}
Réfuter $\sqrt{a+b}=\sqrt a+\sqrt b$ avec des valeurs positives. Démontrer en revanche $\sqrt{ab}=\sqrt a\sqrt b$ pour $a,b\geq0$.
